\scanpage{4}\mspage{1}
% slanted
\margin{I Construction d'un $2$-polyèdre triangulé orienté régulier, associé à un $\Lambda$-module libre de rang $2$, pour $\Lambda\simeq\mathbb{Z}$ ou $\mathbb{Z}/n\mathbb{Z}$\ldots}

Soit $M$ un module libre de rang $2$ sur un anneau $\Lambda$. \deleted{\unsure{Introduisons|Considérons}} Soit $\{x,y,z\}\in M^3$, tel que
\[
\deleted{\unsure{B}}\quad\deleted{B(M)=\{x,\unsure{y}}
\]
\[
x+y+z=0\tag*{(1)}
\]
alors
\[
\{x,y\}\in\deleted{\unsure{\mathrm{Rep}}}\mathrm{Base}_{\Lambda}(M)\iff\{y,z\}\in\deleted{\unsure{\mathrm{Rep}}}\mathrm{Base}_{\Lambda}(M)\iff\{z,x\}\in\unsure{\mathrm{Bas}|\mathrm{Base}}_{\Lambda}(M)\tag*{(2)}
\]
On désigne par
\[
B(M)\subset M^3\tag*{(3)}
\]
l'ens. des triplets $(x,y,z)$ satisfaisant (1) \deleted{\unsure{et}} \deleted{\unsure{qui}} et les conditions équivalentes de (2).

\deleted{Sur} \inserted{La partie} $B(M)$\deleted{, le \unsure{gr.}} de $M^3$ est stable pour l'action de $\mathfrak{S}_3$ opérant sur $M^3$, et $\mathfrak{S}_3$ opère librement sur $B(M)$. Considérons le quotient
\[
\deleted{\unsure{B(M)}}\quad\widetilde F(M)\overset{\text{déf}}{=}B(M)/\mathfrak{S}_3\simeq
\left\{A\subset\unsure{\mathfrak P|\mathscr P}_3(M)\ \middle|\ 
\begin{gathered}\textstyle\sum_{x\in A}x=0,\\ \forall x\in A,\\ A\setminus\{x\}\in\mathrm{Base}\unsure{(M)}\end{gathered}\right\}\tag*{(4)}
\]
\deleted{\unsure{Soit|Notons}} On peut aussi interpréter $B(M)$ comme can. isom. à l'ens. des \uline{repères} de $M$ i.e. des bases ordonnées $(x,y)$, par
\[
\begin{aligned}
B(M)&\xrightarrow{\sim}\mathrm{Rep}(M)\overset{\text{(déf)}}{=}\mathrm{Isom}_{\Lambda}(\Lambda^2,M)\\
(x,y,z)&\longmapsto (x,y)
\end{aligned}\tag*{(5)}
\]
l'action de $\mathfrak{S}_3$ sur $B(M)$ provient alors d'une action de $\mathfrak{S}_3$ sur $\mathrm{Rep}(M)$, déduite
% archive sheet number: 3
